\documentclass[11pt,letterpaper]{article}
\usepackage[letterpaper,textwidth=6.0in,textheight=8.65in,top=1.02in]{geometry}
\usepackage[T1]{fontenc}
\usepackage{lmodern}
\usepackage{amsmath,amssymb,amsthm,booktabs,microtype}
\usepackage[hidelinks,pdfauthor={Blake Stone},pdftitle={Preperiods of modular Ducci sequences}]{hyperref}
\newtheorem{theorem}{Theorem}
\newtheorem{lemma}[theorem]{Lemma}
\newtheorem{proposition}[theorem]{Proposition}
\newcommand{\Z}{\mathbb Z}
\newcommand{\F}{\mathbb F}
\newcommand{\ord}{\operatorname{ord}}
\newcommand{\card}[1]{\lvert #1\rvert}
\newcommand{\ind}{\mathbf1}
\newcommand{\A}{A_{m,n}}
\setcounter{secnumdepth}{1}
\setlength{\parindent}{1.3em}
\raggedbottom
\begin{document}
\begin{center}
{\LARGE\bfseries Preperiods of modular Ducci sequences\par}
\vspace{9pt}
{\large Blake Stone\par}
\vspace{3pt}
{\small November 26, 2026\par}
\end{center}
\begin{abstract}
The modular Ducci map adds adjacent coordinates of a cyclic tuple. We prove
that the upper bound of Lewis and Tefft for its maximum preperiod is attained.
The lower bound follows from the nilpotency index of the augmentation ideal
of a cyclic group ring. We also give an elementary recursion that counts the
tuples of each preperiod for arbitrary width and modulus. It follows that the
distribution within the basin of any cycle depends only on the cycle length.
\end{abstract}

\section{The maximum preperiod}
For $m\ge2$ and $n\ge1$, let $D$ act on $(\Z/m\Z)^n$ by
$(Dv)_i=v_i+v_{i+1}$, with indices read modulo $n$. Write $\tau(v)$ for
the least $h\ge0$ for which $D^{h+c}v=D^hv$ for some $c\ge1$, and set
$L_m(n)=\max_v\tau(v)$. Thus a tuple on a cycle has preperiod zero.
Dular asked for the largest preperiod of this map \cite[\S7]{dular}.
Lewis and Tefft determined it for odd $n$ \cite{odd} and proved a
prime-power upper bound for even $n$. They proposed that equality always
holds \cite[Theorem 1(3) and the paragraph following it]{even}.

For a prime $p$, an integer $\ell\ge1$, and $k\ge0$, put
\begin{equation}\label{H}
 H_p(k,\ell)=
 \begin{cases}
  1,& k=0,\\
  p^{k-1}\bigl(1+\ell(p-1)\bigr),&k\ge1.
 \end{cases}
\end{equation}
Here $v_p(a)$ denotes the exponent of $p$ in $a$, and $p^\ell\parallel m$
means $\ell=v_p(m)>0$.
\begin{theorem}\label{max}
The tuple $(1,0,\dots,0)$ attains the maximum preperiod, which is
\[
 L_m(n)=\begin{cases}
 \displaystyle\max_{p^\ell\parallel m}H_p(v_p(n),\ell),& n\text{ even},\\
 v_2(m),& n\text{ odd}.
 \end{cases}
\]
\end{theorem}
In the case singled out by Lewis and Tefft, $p$ is the only common prime
factor of $n$ and $m$. Writing $k=v_p(n)$ and $\ell=v_p(m)$, their upper
bound is precisely $H_p(k,\ell)$.

\paragraph{The lower bound in one quotient.}
Identify the tuple with a polynomial by numbering in the direction for
which the shift $v_i\mapsto v_{i+1}$ acts as multiplication by $x$.
For even $n$, reduction of coefficients and substitution $x\mapsto-y$ give
a surjection
\[
 (\Z/m\Z)[x]/(x^n-1)\ \longrightarrow\
 (\Z/p^\ell\Z)[y]/(y^{p^k}-1),\qquad k=v_p(n).
\]
The unit seed maps to $1$ and its $h$th iterate maps to $(1-y)^h$.
The cyclic augmentation-ideal theorem of Clark and Schauz
\cite[\S4.3, proof of Theorem 4.7 in the preprint]{cs} gives the exact
nilpotency index $H_p(k,\ell)$ for $k\ge1$. The image first reaches zero at
that time, and zero is its only periodic value. A periodic iterate maps
to a periodic iterate. This proves the conjectured lower bound.

\section{The cyclic components}
Fix $p^\ell\parallel m$ and write $B=\Z/p^\ell\Z$. When $n$ is even,
$x=-y$ identifies $B[x]/(x^n-1)$ with $B[y]/(y^n-1)$. Put
$k=v_p(n)$, $q=p^k$ and $u=n/q$. In the factorization
\[
 y^n-1=(y^q-1)C(y),\qquad C(y)=1+y^q+\cdots+y^{(u-1)q},
\]
the factors are comaximal since $C(y)\equiv u\bmod (y^q-1)$ and $u$ is
a unit. The Chinese remainder theorem splits the ring accordingly.
On the component $B[y]/C$, of size $p^{\ell(n-q)}$, the Ducci multiplier
$1-y$ is invertible since $C(1)=u$. This component is trivial when $u=1$.
On the other component it is $-t$, where
\[
 R_{k,\ell}=B[t]/\bigl((1+t)^{p^k}-1\bigr),\qquad t=y-1.
\]

\begin{lemma}[Cyclic nilpotency index]\label{index}
In $R_{k,\ell}$, the nilpotency index of $t$ is $H_p(k,\ell)$.
\end{lemma}
\begin{proof}
For $k=0$ the element $t$ vanishes. For $k\ge1$, write
$d=p^{k-1}$, $w=q-d$ and $H=q+(\ell-1)w$. The monic defining relation
gives a basis $1,t,\ldots,t^{q-1}$ and a reduction
\[
 t^q=-\sum_{a=1}^{q-1}\binom qa t^a.
\]
Here $v_p\binom qa=k-v_p(a)$, since
$\binom qa=(q/a)\binom{q-1}{a-1}$ and all coefficients of
$(1+t)^{q-1}=(1+t^q)/(1+t)$ in $\F_p[t]$ are units. Assign the
integer monomial $bt^i$ weight $v_p(b)w+i$. Reduction never decreases
weight. Indeed, $\binom qa t^a$ has weight at least $w+a\ge q$ when
$v_p\binom qa=1$, because then $d\mid a$. Equality occurs only at $a=d$.
If the valuation is at least two, its weight is at least $2w+a>q$.

After reduction, a nonzero term modulo $p^\ell$ has weight at most $H-1$.
Thus all paths reducing $t^H$ vanish. Starting with $t^{H-1}$, the unique
path that preserves weight chooses $a=d$ at every reduction. After
$\ell-1$ steps it yields a term of degree $q-1$ and valuation $\ell-1$.
All other paths vanish. Hence $t^{H-1}\ne0$.
\end{proof}

\begin{proof}[Completion of the proof of Theorem~\ref{max}]
An invertible component of a finite ring is periodic under multiplication
by its unit multiplier. A nilpotent component first becomes periodic when
it reaches zero. The splitting above and Lemma~\ref{index} give the
even-width assertion for $B$, with equality on the unit seed.

If $n$ is odd and $p$ is odd, $(1+x)(1-x+\cdots+x^{n-1})=2$ in
$B[x]/(x^n-1)$, so every tuple is periodic. For odd $n$ and $p=2$, factor
$x^n-1=(x-1)C_1(x)$. The factors are comaximal since $C_1(1)=n$ is odd.
On the $C_1$ component, $1+x$ is a unit since $C_1(-1)=1$.
On the $x-1$ component it acts on $B$ by multiplication by two, with index
$\ell$. Finally the Chinese remainder theorem in the coefficient modulus
makes $\tau(v)$ the maximum of the component preperiods. Indeed, after all
components have entered cycles, a common multiple of their periods is a
return time. The unit seed attains all component maxima simultaneously.
\end{proof}

\section{The distribution of preperiods}
Put $\A(h)=\#\{v:\tau(v)\le h\}$ for $h\ge0$ and $\A(-1)=0$.
The difference $\A(h)-\A(h-1)$ counts the tuples of preperiod exactly $h$.
For $s=\min(k,\ell)$ define
\begin{equation}\label{F}
 F_{p,k,\ell}(h)=(\ell-s)\min(h,1)+
 \sum_{i=1}^{s}\min\bigl(h,H_p(k-i+1,\ell-i+1)\bigr).
\end{equation}
\begin{theorem}\label{dist}
For every integer $h\ge0$,
\[
 \A(h)=
 \begin{cases}
 \displaystyle\prod_{p^\ell\parallel m}
 p^{\ell(n-p^{k})+F_{p,k,\ell}(h)},\quad k=v_p(n)\text{ in each factor},
       & n\text{ even},\\[3pt]
 m^n\,2^{\min(h,b)-b},\quad b=v_2(m),& n\text{ odd}.
 \end{cases}
\]
\end{theorem}
On the odd-width component where $D$ multiplies by two, exactly
$2^{\min(\ell,h)}$ elements are killed by step $h$. All invertible
components are unrestricted. This proves the odd case. The even case
reduces to showing that $F_{p,k,\ell}(h)$ equals
$\log_p\card{\ker(t^h\mid R_{k,\ell})}$.

\paragraph{The successive quotients.}
Fix $p$, write $H(k,\ell)=H_p(k,\ell)$, and put $R=R_{k,\ell}$. For $j\ge0$
let $e_{k,\ell}(j)=\log_p\card{t^jR/t^{j+1}R}$. The quotient is cyclic
as a $B$-module, generated by $t^j$. Write $\nu_\ell$ for the $p$-valuation
truncated at $\ell$, so that $\nu_\ell(0)=\ell$. Then
\begin{equation}\label{ann}
 e_{k,\ell}(j)=\min\{\nu_\ell(V(0)): V\in R,\ t^jV=0\}.
\end{equation}
In fact $at^j\in t^{j+1}R$ if and only if $a$ is the constant coefficient
of an annihilator of $t^j$. Evaluation at $t=0$ is well defined. For
$j<H(k,\ell)$ the right side is at least one, since $t$ is nilpotent and
an element with unit constant coefficient is a unit.

\begin{proposition}\label{rec}
For $k\ge1$, $\ell\ge2$ and every $j\ge0$,
\[
 e_{k,\ell}(j)=\ind_{\{j<H(k,\ell)\}}+e_{k-1,\ell-1}(j).
\]
\end{proposition}
\begin{proof}
Put $d=p^{k-1}$, $w=(p-1)d$ and $H'=H(k-1,\ell-1)$. The monic factor
\[
 g(t)=\frac{(1+t)^{p^k}-1}{(1+t)^d-1}=\sum_{a=0}^{p-1}(1+t)^{da}
\]
has degree $w$, constant coefficient $p$ as an integer polynomial, and
reduces to $t^w$ modulo $p$. Set $S=B[t]/g$ and let $\pi$ be the image of
$t$. Writing $g=t^w+pU$ with $U(0)=1$ gives $\pi^{\ell w}=0$.
The quotient $S/(\pi)$ is $\F_p$. Every $(\pi^i)/(\pi^{i+1})$ has at most
$p$ elements. Since $\card{S}=p^{\ell w}$ by monicity, all $\ell w$
successive quotients down to zero have exactly $p$ elements. Every nonzero
element of $S$ thus equals a unit times $\pi^i$ for a unique $0\le i<\ell w$:
factor out the largest power of $\pi$, and observe that any element outside
$(\pi)$ is a unit, since $(\pi)$ is nilpotent and its residue ring is a
field. It follows that $\operatorname{Ann}_S(\pi^j)=(\pi^{\ell w-j})$
for $0\le j\le\ell w$. Also $U(\pi)$ is a unit, so $(p)=(\pi^w)$ in $S$.

Consider $0\le j<H'$ and let $V\in R$ annihilate $t^j$. The smaller
ring $R_{k-1,\ell-1}$ has size $p^{(\ell-1)d}$ and its powers of $t$ descend
strictly until zero. Hence $H'\le(\ell-1)d$. In particular $\ell w-j\ge d$.
Under the quotient map $R\to S$, the image $V(\pi)$ lies in
$(\pi^{\ell w-j})\subseteq(\pi^d)$. Because $d\le w$,
\[
 S/(\pi^d)\cong\F_p[t]/t^d=R_{k-1,1}.
\]
Consequently the image of $V$ under $R\to R_{k-1,\ell}$ is $pG$ for some
$G$. The defining polynomial in this smaller ring is monic, so its
underlying $B$-module is free. From $pt^jG=0$ we obtain $t^jG=0$ modulo
$p^{\ell-1}$. Evaluation at zero gives $V(0)=pG(0)$ in $B$.
Equation~\eqref{ann} now yields $e_{k,\ell}(j)\ge1+e_{k-1,\ell-1}(j)$.

Write $H=H(k,\ell)$. We have $H'\le H$, for $R_{k-1,\ell-1}$ is a quotient
of $R$. The smaller exponents vanish for $j\ge H'$, while the larger ones
are at least one for $j<H$. The claimed equality therefore holds as a
lower bound for every $j$. Sum both sides over all $j$. The left side sums
to $\ell p^k$, and the right to $H+(\ell-1)d=\ell p^k$ by~\eqref{H}.
Every term must agree.
\end{proof}

\begin{proof}[Completion of the proof of Theorem~\ref{dist}]
A finite abelian group endomorphism has equally large kernel and cokernel. Thus
$\log_p\card{\ker(t^h\mid R)}=\log_p\card{R/t^hR}=\sum_{j=0}^{h-1}e_{k,\ell}(j)$.
Summing Proposition~\ref{rec} gives the recursion
\[
 F_{p,k,\ell}(h)=\min(h,H(k,\ell))+F_{p,k-1,\ell-1}(h).
\]
When $\ell=1$, $R_{k,1}=\F_p[t]/t^{p^k}$ yields
$\min(h,p^k)$. When $k=0$, multiplication by $t$ is zero and the exponent
is $\ell\min(h,1)$. Iteration gives~\eqref{F}. Multiplying by the sizes of
the invertible components completes the count.
\end{proof}

\paragraph{A four-coordinate example.}
For $n=4$, $m=8$, the local exponent is
$\min(h,1)+\min(h,3)+\min(h,8)$. All tuples eventually reach zero. The
individual counts illustrate the change in growth at steps $1$, $3$ and $8$.
\begin{center}
\begin{tabular}{@{}lrrrrrrrrr@{}}
\toprule
Preperiod $h$ & 0&1&2&3&4&5&6&7&8\\
Tuples & 1&7&24&96&128&256&512&1024&2048\\
\bottomrule
\end{tabular}
\end{center}
For a composite modulus, the maxima combine by taking their maximum.
For example, $L_{72}(12)=\max(H_2(2,3),H_3(1,2))=\max(8,5)=8$.

\paragraph{The basin of a specified cycle.}
The invertible component consists of $\A(0)$ periodic states. A cycle of
length $c$ has $c$ such states in its basin. Every periodic state can be
paired with the same nilpotent states of a given preperiod, so the number
in that basin with preperiod exactly $h$ is
\[
 \frac{c}{\A(0)}\bigl(\A(h)-\A(h-1)\bigr).
\]

\paragraph{How often is the maximum attained?}
At least half the tuples have preperiod $L=L_m(n)$. For $L=0$ all do.
For $L>0$, the tuples of preperiod at most $L-1$ form a proper additive
subgroup: their nilpotent components belong to the kernels of
multiplication by the $(L-1)$st power of the local multiplier. The index
of this subgroup is at least two. More precisely, when $n$ is even,
$m=p^\ell$, and $k=v_p(n)>0$, the proportion of tuples with the
maximum preperiod is $1-1/p$. Indeed $H_p(k,\ell)>1$ is the unique
largest cutoff in~\eqref{F}, so $F_{p,k,\ell}(L)-F_{p,k,\ell}(L-1)=1$.

\paragraph{A coordinate test for periodicity.}
For even $n$ and each $p^\ell\parallel m$, put $q=p^{v_p(n)}$. A tuple is
periodic precisely when its projection on every nilpotent component is zero.
In the polynomial coordinates used above this is the system
\[
 \sum_{\substack{0\le j<n\\ j\equiv a\pmod q}}(-1)^jv_j
 \equiv0\pmod{p^\ell}\qquad(0\le a<q).
\]
These are exactly the coefficients obtained by reducing the tuple
polynomial at $x=-y$ modulo $y^q-1$. Reversing the cyclic numbering
permutes the congruences since $n$ is even. For odd $n$ the only condition
is $\sum_jv_j\equiv0\bmod 2^{v_2(m)}$, in agreement with \cite{odd}.

\paragraph{Relation to the lens-space computation.}
The integral cokernel underlying the local count at $h\ge1$ is
$\Z[t]/(t^h,(1+t)^{p^k}-1)$. This is also the complex $K^0$-ring of the
standard lens space $S^{2h-1}/C_{p^k}$. Indeed this space is the circle
bundle of the $p^k$th tensor power of the tautological line on
$\mathbb{CP}^{h-1}$, up to dualizing. Its $K$-theory Gysin sequence gives
exactly the displayed quotient of $K^0(\mathbb{CP}^{h-1})=\Z[t]/t^h$.
The additive decomposition of this group was computed by Mahammed and
Yanagida \cite{ma,ya}. In particular it determines the same kernel sizes
after reduction modulo $p^\ell$. Proposition~\ref{rec} computes these
sizes directly with finite rings.

\paragraph{Acknowledgment.}
The author used OpenAI Codex in developing the proofs, checking examples,
and preparing the manuscript.

\begin{thebibliography}{9}\small
\bibitem{cs} P.~L. Clark and U. Schauz, Functional degrees and arithmetic
applications I: the set of functional degrees, \emph{J. Algebra}
\textbf{608} (2022), 691--718. \href{https://arxiv.org/abs/2201.02763}{arXiv:2201.02763}.
\bibitem{dular} B. Dular, Cycles of sums of integers, \emph{Fibonacci Quart.}
\textbf{58} (2020), 126--139. \href{https://arxiv.org/abs/1905.01765}{arXiv:1905.01765}.
\bibitem{odd} M.~L. Lewis and S.~M. Tefft, Ducci on $\Z_m^n$ and the maximum
length for $n$ odd (2024). \href{https://arxiv.org/abs/2403.05319}{arXiv:2403.05319}.
\bibitem{even} M.~L. Lewis and S.~M. Tefft, The maximum length for Ducci
sequences on $\Z_m^n$ when $n$ is even, \emph{Integers} \textbf{26} (2026), A17.
\href{https://math.colgate.edu/~integers/aa17/aa17.pdf}{Journal article}.
\bibitem{ma} N. Mahammed, $K$-th\'eorie des espaces lenticulaires,
\emph{C. R. Acad. Sci. Paris S\'er. A} \textbf{272} (1971), 1363--1365.
\bibitem{ya} S. Yanagida, The additive structure of $G^*(L^n(p^k))$,
\emph{Osaka J. Math.} \textbf{13} (1976), 523--536.
\href{https://doi.org/10.18910/11698}{doi:10.18910/11698}.
\end{thebibliography}
\end{document}
